\section{The bibliography of EXI-01-a}\label{s:bib}

The bibliography of EXI-01-a has three items, one for each of the three positions that the text must locate itself among: \citet{Wolf2010} for the hybrid view it adopts, \citet{Metz2013} for the objectivism about worth on which it relies, and \citet{Nagel1971} for the challenge it does not answer. The selection is a summary of the text's dialectical situation. Each item is cited once, at the point where the position it represents enters the argument: Wolf at the definition, Metz at the assumption of objective worth, Nagel at the remark on what the theorem leaves open.

\paragraph{Wolf.} The citation is of the Tanner Lectures, not of the earlier paper \citep{Wolf1997} in which the fitting-fulfilment view was first stated. The lectures are the fuller statement, they include the defence against subjectivism on which Definition~1 relies, and they are the version most often cited in the discipline: EXI-17-a, EXI-20-a, EXI-01-c and EXI-01-d all cite them. The Commentary has cited the earlier paper where its formulations differ (\Lm{2.8}).

\paragraph{Metz.} The citation is of Metz's book, which is both the standard survey of theories of meaning and the statement of his own theory. The comment on \Lm{3.5} has noted that Metz is an ally of EXI-01-a on the existence of objective worth and a rival on the measure of meaning. A reader who wants Metz's survey of the debate will also find his classification of the views from which EXI-01-a distinguishes itself, and the Commentary's introduction follows that classification.

\paragraph{Nagel.} The citation is of the 1971 paper, not of \emph{The View from Nowhere}. The paper is where the sense of the absurd is analyzed as a collision between seriousness and detachment and where the ironic response is proposed, and it is the paper that EXI-04-a reconstructs. The book generalizes the analysis and adds a treatment of death; the Commentary cites it at \Lm{11.2} and Volume~II examines the relation between the two.

\paragraph{What is not cited.} Four works that the Commentary has found indispensable for understanding EXI-01-a are absent from its bibliography. Taylor's chapter on the meaning of life \citep[ch.~18]{Taylor1970} is the source of the subjectivist view that the qualifier \enquote{of objective worth} excludes, and of the figure of Sisyphus that the discipline uses throughout; EXI-17-a cites it. Nozick's chapter \citep[ch.~6]{Nozick1981} is the principal statement of the view that meaning requires connection with value beyond oneself; Definition~1 secures connection through objective worth, and EXI-17-a's reply to its Objection~2 explains why connection does not require success. Camus's essay \citep{Camus1955} and Tolstoy's \emph{Confession} \citep{Tolstoy1987} are the classical texts of the absurd and of the crisis of meaning. Their absence from a one-page text is not a judgment of their importance, and the Commentary has supplied them where they bear on the lemmata.
